%! Multiplying \& Dividing Radicals; Rationalizing — Unit 5, Lesson 5.3 — Guided Notes
\documentclass[10pt]{article}
\usepackage{algebra2-article}
\usepackage{algebra2-boxes}
\newcommand{\TallMath}[1]{$\displaystyle #1\rule[-1.4em]{0pt}{3.2em}$}

\begin{document}

\pageheader{Unit 5, Lesson 5.3}{Guided Notes}
\namedateperiod

\vspace{0.08in}

\begin{objectivebox}
\textbf{By the end of this lesson, I will be able to\dots}
\begin{itemize}\setlength{\itemsep}{2pt}
  \item \emph{multiply} radical expressions --- including distributing over a sum and multiplying two binomials --- and simplify the result;
  \item recognize a \emph{conjugate} and explain why a conjugate pair always multiplies to a \emph{rational} number;
  \item \emph{rationalize} a denominator, with one radical or with a binomial, and explain why the value of the expression never changes.
\end{itemize}
\end{objectivebox}

\vspace{0.05in}

\begin{vocabbox}
\small
Fill in each term as we build it together.
\par\vspace{2pt}   % \par is required: \termblanklong's \noindent is a no-op mid-paragraph
\termblanklong{Rationalizing the denominator}
\termblanklong{Conjugate}
\termblanklong{Multiplying by a form of $1$}
\termblanklong{Simplest radical form (all three conditions, at last)}
\end{vocabbox}

\begin{hookbox}
Two students are asked for the exact value of $\dfrac{1}{\sqrt{2}}$. One writes $\dfrac{1}{\sqrt2}$
and stops. The other writes $\dfrac{\sqrt2}{2}$. Their teacher marks \emph{both} correct.
\begin{center}
\renewcommand{\arraystretch}{1.6}
\begin{tabularx}{\linewidth}{@{}X X@{}}
$\dfrac{1}{\sqrt2} \approx \dfrac{1}{1.414\,214}=$ \blank{2.4cm} &
$\dfrac{\sqrt2}{2} \approx \dfrac{1.414\,214}{2}=$ \blank{2.4cm} \\
\end{tabularx}
\end{center}
They are the same number. So here is the real question, and it is a fair one:
\textbf{if both are correct, why does every textbook insist on the second one?}
\vspace{2pt}

Imagine it is $1850$ and you must actually \emph{compute} one of these by hand, with no calculator.
Which long division would you rather do --- \; $1 \div 1.414\,214$ \; or \;
$1.414\,214 \div 2$? \writeline
\vspace{2pt}
That is the whole reason the rule exists. A radical on \emph{top} is harmless; a radical on the
\emph{bottom} is a nuisance. Now look at how the second student got there: they multiplied top and
bottom by $\sqrt2$, and $\sqrt2\cdot\sqrt2=\blank{0.9cm}$ --- the radical simply
\blank{2.0cm}. That one fact runs today's entire lesson.
\end{hookbox}

\begin{notesbox}{1. Multiplying: the Product Property, Run Forwards}
Yesterday you read $\sqrt[n]{ab}=\sqrt[n]{a}\cdot\sqrt[n]{b}$ from \emph{left to right} to pull
factors out. Today you read the same statement from \emph{right to left}.
\begin{center}
\begin{tcolorbox}[colback=goldbg, colframe=goldacc, boxrule=0.8pt, arc=1.5mm,
    left=3mm, right=3mm, top=1.5mm, bottom=1.5mm, width=0.92\linewidth]
\centering
$\displaystyle \sqrt[n]{a}\cdot\sqrt[n]{b}=\sqrt[n]{ab}$ \qquad --- \qquad
multiply the \textbf{outsides} together, multiply the \textbf{insides} together, \emph{then}
simplify.
\end{tcolorbox}
\end{center}

\vspace{2pt}
\textbf{Simplify \emph{last}, and never skip it.} A product often contains a perfect power that
neither factor had on its own:
\[
  \sqrt{6}\cdot\sqrt{10}=\sqrt{\blank{1.0cm}}=\blank{1.8cm}
  \qquad\qquad
  \sqrt{3}\cdot\sqrt{12}=\sqrt{\blank{1.0cm}}=\blank{1.4cm}
\]
Neither $\sqrt6$ nor $\sqrt{10}$ could be simplified alone. Multiplying them created a perfect
square out of nothing --- so \emph{finishing} is not optional.

\vspace{3pt}
\textbf{With coefficients out front.} Outsides multiply outsides; insides multiply insides.
\begin{center}
\renewcommand{\arraystretch}{1.9}
\begin{tabularx}{\linewidth}{@{}X X X@{}}
(a) \; $2\sqrt5\cdot3\sqrt{10}=$ \blank{2.2cm} & (b) \; $\sqrt[3]{4}\cdot\sqrt[3]{6}=$ \blank{2.2cm} &
(c) \; $\sqrt{2x}\cdot\sqrt{6x^{3}}=$ \blank{2.4cm} \\
\end{tabularx}
\end{center}

\vspace{2pt}
\textbf{The one fact the whole lesson rests on.} A square root multiplied by itself:
\[
  \sqrt{a}\cdot\sqrt{a}=\sqrt{a^{2}}=\blank{1.0cm}
  \qquad\text{so}\qquad \left(\sqrt{7}\right)^{2}=\blank{1.0cm},
  \qquad \left(3\sqrt5\right)^{2}=9\cdot 5=\blank{1.0cm}
\]
\textbf{A square root times itself erases the radical.} (For a cube root you need
\blank{1.2cm} copies: $\sqrt[3]{2}\cdot\sqrt[3]{2}\cdot\sqrt[3]{2}=2$. The index still sets the
group size.)

\vspace{3pt}
\emph{One restriction.} The indices must \blank{1.6cm} before you may combine two radicals. For
$\sqrt{2}\cdot\sqrt[3]{2}$ there is no radical rule at all --- but Lesson 5.1 handles it:
$2^{1/2}\cdot 2^{1/3}=2^{\,\blank{1.0cm}}=\sqrt[6]{\blank{0.8cm}}$.
\end{notesbox}

\begin{notesbox}{2. Distributing and FOIL --- Radicals Behave Like Binomials}
Nothing new here at all. In Warm-Up item~1 you expanded $3x(2x+5)$ and $(x+4)(x-4)$. Radicals
follow the identical rules; the radical just plays the role of the variable.
\[
  \sqrt{5}\left(\sqrt{3}+\sqrt{10}\right)
  = \sqrt{\blank{0.8cm}}+\sqrt{\blank{0.8cm}} = \blank{2.6cm}
  \qquad
  \sqrt{2}\left(4-\sqrt{2}\right)=\blank{2.6cm}
\]
\emph{Simplify each product separately after distributing} --- in the first one, only the second
term simplifies.

\vspace{3pt}
\textbf{Two binomials: FOIL, then collect like radicals (Lesson 5.2).}
\[
  \left(2+\sqrt3\right)\left(5-\sqrt3\right)
  = 10 \;\blank{1.2cm}\; 2\sqrt3 \;+\; 5\sqrt3 \;\blank{1.2cm}\; 3
  = \blank{2.6cm}
\]

\vspace{3pt}
\textbf{Squaring a binomial --- the error of the day.}
\[
  \left(\sqrt7+\sqrt2\right)^{2}
  = \left(\sqrt7+\sqrt2\right)\left(\sqrt7+\sqrt2\right)
  = 7 + \blank{1.4cm} + \blank{1.4cm} + 2 = \blank{2.6cm}
\]
It is \textbf{not} $7+2=9$. There is a middle term, for exactly the reason
$(x+4)^{2}\ne x^{2}+16$ --- which you already rejected in the Warm-Up. \; Your turn:
\; $\left(3-\sqrt5\right)^{2}=\blank{2.8cm}$
\end{notesbox}

\begin{notesbox}{3. The Pair That Comes Out Rational: Conjugates}
Now multiply a binomial by its \emph{sign-flipped} partner --- Warm-Up item 1(b), with a radical in
place of the $4$.
\[
  \left(4+\sqrt3\right)\left(4-\sqrt3\right)
  = 16 - 4\sqrt3 + 4\sqrt3 - 3 = \blank{1.4cm}
\]
Two things happened at once, and both matter. The middle terms \blank{2.0cm}, and the last term
lost its radical because $\sqrt3\cdot\sqrt3=\blank{0.8cm}$. What is left is a plain
\blank{2.2cm} number --- no radical anywhere.
\begin{center}
\begin{tcolorbox}[colback=goldbg, colframe=goldacc, boxrule=0.8pt, arc=1.5mm,
    left=3mm, right=3mm, top=1.5mm, bottom=1.5mm, width=0.92\linewidth]
\centering
\textbf{Conjugates:} $a+\sqrt{b}$ and $a-\sqrt{b}$ --- same two terms, \textbf{opposite} middle
sign. \\[3pt]
$\displaystyle \left(a+\sqrt b\right)\left(a-\sqrt b\right)=a^{2}-b$
\qquad
$\displaystyle \left(\sqrt a+\sqrt b\right)\left(\sqrt a-\sqrt b\right)=a-b$ \\[3pt]
{\small Both products are \textbf{always rational}. It is just difference of squares, and squaring
is what undoes a square root.}
\end{tcolorbox}
\end{center}

\vspace{2pt}
\textbf{You have seen this move before.} In Lesson 2.4 you multiplied $a+bi$ by $a-bi$ and got
$a^{2}+b^{2}$ --- a real number. Same trick, different number system: there it turned an
\emph{imaginary} number real; here it turns an \emph{irrational} number rational. \; Why the
different sign? Because $i^{2}=\blank{1.0cm}$ while $\left(\sqrt b\right)^{2}=\blank{1.0cm}$.

\vspace{3pt}
\textbf{Practice.} Write each conjugate, then multiply the pair.
\begin{center}
\renewcommand{\arraystretch}{1.7}
\begin{tabularx}{\linewidth}{@{}l X X@{}}
\hline
\textbf{Expression} & \textbf{Its conjugate} & \textbf{Their product} \\
\hline
$6+\sqrt{5}$          & \blank{2.6cm} & \blank{2.2cm} \\
$\sqrt{7}-2$          & \blank{2.6cm} & \blank{2.2cm} \\
$\sqrt{11}+\sqrt{3}$  & \blank{2.6cm} & \blank{2.2cm} \\
\hline
\end{tabularx}
\end{center}
\end{notesbox}

\begin{notesbox}{4. Dividing: the Quotient Property, Run Forwards}
Same idea, one floor down: two radicals with the \emph{same index} can be pulled under one radical.
\begin{center}
\renewcommand{\arraystretch}{2.0}
\begin{tabularx}{\linewidth}{@{}X X X@{}}
(a) \; $\dfrac{\sqrt{48}}{\sqrt{3}}=$ \blank{1.8cm} &
(b) \; $\dfrac{6\sqrt{20}}{2\sqrt{5}}=$ \blank{1.8cm} &
(c) \; $\dfrac{\sqrt{40x^{3}}}{\sqrt{5x}}=$ \blank{2.2cm} \\
\end{tabularx}
\end{center}
Every one of those came out clean. But this one does not:
\; $\dfrac{\sqrt3}{\sqrt5}=\sqrt{\dfrac{3}{5}}$ --- and now there is a radical sitting in the
\blank{2.2cm}, which breaks simplest-form condition \blank{0.8cm}. That is the condition we
skipped yesterday. Time to finish it.
\end{notesbox}

\begin{notesbox}{5. Rationalizing a Denominator with One Radical}
\textbf{The move:} multiply the top \emph{and} the bottom by the radical that is stuck in the
bottom. From Warm-Up item~3, that is legal because $\dfrac{\sqrt b}{\sqrt b}=\blank{0.8cm}$, and
multiplying by $1$ cannot change a number's value --- only its \blank{1.6cm}.
\[
  \frac{1}{\sqrt2}\cdot\frac{\sqrt2}{\sqrt2}
  = \frac{\sqrt2}{\blank{0.8cm}}
  \qquad\qquad
  \frac{3}{\sqrt5}\cdot\frac{\blank{1.0cm}}{\blank{1.0cm}} = \blank{1.8cm}
\]
The Hook is now closed: the two students had the same number the whole time.

\vspace{3pt}
\textbf{Your turn.} Rationalize each. In (c), \emph{simplify the denominator first}.
\begin{center}
\renewcommand{\arraystretch}{2.2}
\begin{tabularx}{\linewidth}{@{}X X X@{}}
(a) \; $\dfrac{5}{\sqrt{3}}=$ \blank{1.8cm} &
(b) \; $\sqrt{\dfrac{3}{5}}=$ \blank{1.8cm} &
(c) \; $\dfrac{\sqrt{7}}{\sqrt{12}}=$ \blank{2.0cm} \\
\end{tabularx}
\end{center}

\vspace{2pt}
\textbf{A higher index changes what you multiply by.} For $\dfrac{1}{\sqrt[3]{2}}$, multiplying by
$\sqrt[3]{2}$ is \emph{not} enough --- it only gives $\sqrt[3]{4}$ on the bottom, and that is still
a radical. You need \blank{1.2cm} equal factors to escape a cube root, so multiply by
$\sqrt[3]{4}$:
\[
  \frac{1}{\sqrt[3]{2}}\cdot\frac{\sqrt[3]{4}}{\sqrt[3]{4}}
  = \frac{\sqrt[3]{4}}{\sqrt[3]{\blank{0.8cm}}} = \frac{\sqrt[3]{4}}{\blank{0.8cm}}
\]
This is yesterday's \emph{``the index sets the group size''} showing up again: you multiply by
whatever \textbf{completes the group}.
\end{notesbox}

\begin{notesbox}{6. Rationalizing a Binomial Denominator: Use the Conjugate}
If the bottom is $3-\sqrt5$, multiplying by $\sqrt5$ is useless --- try it and you still have a
radical down there. Multiply by the \textbf{conjugate} instead, because a conjugate pair is
guaranteed to come out rational (Section 3).
\[
  \frac{6}{3-\sqrt5}\cdot\frac{\blank{1.6cm}}{\blank{1.6cm}}
  = \frac{6\left(3+\sqrt5\right)}{9-\blank{0.8cm}}
  = \frac{6\left(3+\sqrt5\right)}{\blank{0.8cm}}
  = \blank{2.6cm}
\]
Notice the last step: the numerator and the denominator shared a common factor, so the answer
reduced. \emph{Always check whether it does.}

\vspace{3pt}
\textbf{Your turn.}
\begin{center}
\renewcommand{\arraystretch}{2.4}
\begin{tabularx}{\linewidth}{@{}X X@{}}
(a) \; $\dfrac{4}{\sqrt5-1}=$ \blank{2.8cm} &
(b) \; $\dfrac{\sqrt2}{\sqrt7+\sqrt3}=$ \blank{3.0cm} \\
\end{tabularx}
\end{center}

\vspace{2pt}
\textbf{Three ways this goes wrong.} Write the fix next to each.
\begin{itemize}[leftmargin=1.5em, itemsep=3pt]
  \item Multiplying only the \emph{denominator} by the conjugate. \; Why that is illegal:
        \blank{5.0cm}
  \item Multiplying by the \emph{same} binomial instead of the conjugate. \;
        $\left(3-\sqrt5\right)^{2}=\blank{2.2cm}$, which is still \blank{1.8cm}.
  \item Forgetting to FOIL the \emph{numerator} when it is also a binomial. \;
        $\left(2+\sqrt3\right)\left(1+\sqrt3\right)$ has \blank{1.0cm} products, not $2$.
\end{itemize}
\end{notesbox}

\begin{practicebox}
Work these with your class. Assume all variables represent positive real numbers, and leave no
radical in any denominator.
\begin{enumerate}[leftmargin=1.7em, itemsep=5pt]
  \item Multiply: \; $\sqrt{6}\cdot\sqrt{15}=\blank{1.8cm}$ \quad
        $3\sqrt2\cdot4\sqrt{6}=\blank{1.8cm}$ \quad $\sqrt[3]{9}\cdot\sqrt[3]{6}=\blank{1.8cm}$
  \item Expand: \; $\sqrt3\left(\sqrt{12}-\sqrt6\right)=\blank{2.2cm}$ \quad
        $\left(4+\sqrt2\right)\left(3-\sqrt2\right)=\blank{2.2cm}$
  \item Expand: \; $\left(\sqrt5+\sqrt3\right)^{2}=\blank{2.4cm}$ \quad
        $\left(2\sqrt7-3\right)\left(2\sqrt7+3\right)=\blank{2.0cm}$
  \item Rationalize: \; $\dfrac{8}{\sqrt6}=\blank{1.8cm}$ \quad
        $\dfrac{3}{5+\sqrt2}=\blank{2.4cm}$
  \item \textbf{Explain.} Both answers in item 3 came from squaring or multiplying a pair of
        binomials with the same two terms, yet one answer still has a radical in it and the other
        does not. What is the difference between the two problems? \writelines{2}
\end{enumerate}
\end{practicebox}

\end{document}
